\section{Elevator pitchs y preguntas frecuentes}

En esta sección se documentan las distintas variantes de presentación rápida (Elevator Pitch) diseñadas para diferentes audiencias, así como las respuestas estructuradas para los roles clave del equipo ante consultas de inversores, tutores o clientes.

\subsection{Variantes de elevator pitch}

\subsubsection{Variante 1: Enfoque comercial e inversores}

¿Sabían que las tres principales fugas de capital en los clubes deportivos provienen de la morosidad no detectada, los ingresos no autorizados en los estadios y la falta de optimización en el alquiler de espacios compartidos?

Soy [Nombre], parte del equipo creador de SocioUnido. Para resolver este problema de raíz, desarrollamos la única herramienta del mercado que ofrece una digitalización verdaderamente integral del club.

Nuestra plataforma inteligente centraliza absolutamente toda la gestión: desde el cobro automatizado de cuotas y el control seguro de accesos mediante códigos QR, hasta la administración detallada de disciplinas deportivas y la reserva de instalaciones.

Facilitamos y modernizamos los procesos en todos los niveles, conectando a la administración, los empleados y los usuarios en un mismo ecosistema.

Nuestro lema es claro: 'Porque el club es de los socios, y la gestión es de Socio Unido'.

Hoy estamos abriendo nuestra primera ronda de capital para escalar esta solución a nivel regional, y buscamos socios estratégicos que quieran ser parte de la transformación digital del deporte.

Dicho esto queremos preguntarles ¿Quién está listo para ser parte del futuro del fútbol a nivel global?

\subsubsection{Variante 2: Enfoque técnico y académico}

Hola, soy [Nombre] y en representación de nuestro equipo les presento SocioUnido, la única plataforma integral concebida para la digitalización absoluta de instituciones deportivas.

Si bien es de público conocimiento que los clubes pierden millones anualmente debido a la morosidad y a la ineficiente gestión de sus espacios, el verdadero desafío de ingeniería detrás de esta transformación radica en la infraestructura: la alta concurrencia masiva y la nula conectividad en las puertas de los estadios.

Para resolver la gestión diaria, construimos un ecosistema digital —compuesto por un panel web y una app móvil— soportado íntegramente por una arquitectura de microservicios en la nube, lo que nos permite escalar de forma elástica e independiente ante picos de demanda.

Sin embargo, nuestra mayor innovación tecnológica se encuentra en el control de accesos. Hemos desarrollado un sistema de validación de entradas 100\% offline. Mediante el uso de criptografía y la generación dinámica de códigos, garantizamos un ingreso seguro, rápido y tolerante a caídas de red, eliminando por completo el histórico cuello de botella en los molinetes.

A continuación, los invitamos a sumergirse en la arquitectura y el funcionamiento de nuestra plataforma.

\newpage

\subsubsection{Variante 3: Enfoque a clubes / dirigentes (Pitch de ventas)}

¿Sabían que, más allá de los trámites lentos que muchas veces alejan a los socios, su institución está sufriendo una fuga de capital por tres ineficiencias silenciosas? La morosidad no detectada, los ingresos no autorizados en los eventos y la subutilización de los espacios compartidos.

Soy [Nombre], y vengo a presentarles SocioUnido. Nosotros creamos una plataforma que pone el club en el bolsillo del socio y el control financiero total en las manos de la dirigencia. 

Para resolver estos problemas de raíz, nuestra aplicación automatiza el cobro de cuotas y el seguimiento de pagos para erradicar la morosidad; centraliza la reserva de canchas y disciplinas maximizando su ocupación; y blinda sus instalaciones con nuestra mayor innovación tecnológica: un sistema de control de accesos por QR 100\% offline, garantizando un ingreso ágil, masivo y antifraude, sin depender de la conexión a internet.

SocioUnido es una solución integral. Su diseño intuitivo asegura una adopción fácil e inmediata, tanto para el personal administrativo como para los socios de todas las edades, sin necesidad de adquirir hardware costoso.

¿Están listos para frenar la fuga de ingresos, modernizar sus procesos y ofrecerle a su comunidad la experiencia de primer nivel que se merecen?

\subsubsection{Variante 4: Enfoque networking y contacto directo}

Mucho gusto [Nombre del cliente], soy [Nombre] de SocioUnido y te quiero presentar nuestra solución.

Nuestra plataforma es exactamente el sistema de digitalización integral que su club necesita hoy para dar un salto definitivo en su gestión.

Sabemos que sus principales puntos de pérdida de dinero hoy son tres: la morosidad acumulada por falta de seguimiento, los ingresos no autorizados en los estadios y la subutilización de los espacios compartidos.

Nuestro sistema lo resuelve de raíz centralizando toda la gestión.

Le damos a la administración un panel de control total potenciado con inteligencia artificial para detectar patrones de morosidad y recuperar socios deudores, al socio una aplicación móvil para que pague sus cuotas y reserve espacios en segundos. Y además, blindamos las puertas de su club con un sistema de acceso por código QR 100\% offline, que garantiza un ingreso rápido y antifraude sin depender de internet.

La digitalización y el uso de nuestra plataforma es extremadamente simple. SocioUnido es una solución integral, con un diseño muy intuitivo tanto para empleados como socios que no requiere hardware adicional.

Me gustaría dejarte mi tarjeta y que nos podamos tomar un café pronto para poder contarte más acerca de cómo podemos llevarlos al siguiente nivel de la gestión deportiva.

\newpage

\subsection{Preguntas frecuentes generales}

\subsubsection{Felipe Ascencio - Product Owner y Lógica de Pagos}

\textbf{P: ¿Cuál es tu rol en el equipo?} \\
\textbf{R:} Como Product Owner, mi responsabilidad es maximizar el valor del producto, priorizando el Backlog y asegurando que las historias de usuario estén alineadas con las necesidades reales de los clubes. Además, lidero la integración de la lógica de pagos, un módulo crítico para solucionar el problema de la morosidad.

\textbf{P: ¿Cuál es la visión a futuro del producto?} \\
\textbf{R:} La visión es convertir a SocioUnido en el estándar de gestión deportiva en la región, logrando que los clubes dependan de una única plataforma centralizada y abandonen los sistemas fragmentados o el papel.

\textbf{P: ¿Qué necesitan actualmente para avanzar?} \\
\textbf{R:} Actualmente buscamos validar nuestra pasarela de pagos en un entorno real (beta testing) con un club de prueba para ajustar las métricas de retención y asegurar una tasa de éxito del 99.9\% en las transacciones.

\subsubsection{Lautaro Ghosn - Infraestructura Cloud, DevOps y Arquitectura}

\textbf{P: ¿Cuál es tu rol en el equipo?} \\
\textbf{R:} Soy el responsable de la columna vertebral del sistema. Diseño la arquitectura de microservicios, gestiono el API Gateway y me encargo de la cultura DevOps, asegurando que los despliegues a la nube (CI/CD) sean automáticos, seguros y tolerantes a fallos.

\textbf{P: ¿Por qué eligieron microservicios en lugar de un monolito?} \\
\textbf{R:} Porque necesitamos escalar de forma elástica. Si el día de un partido hay un pico de tráfico en el control de accesos, nuestra arquitectura permite asignarle más recursos solo a ese microservicio sin sobrecargar el panel administrativo o el módulo de reservas.

\subsubsection{Martín Guerrero - Seguridad, Base de Datos e IA Conversacional}

\textbf{P: ¿Cuál es tu rol en el equipo?} \\
\textbf{R:} Me encargo de garantizar la integridad de la plataforma. Implemento los sistemas de autenticación y autorización (RBAC), administro las bases de datos y desarrollo el módulo de Inteligencia Artificial Conversacional para soporte automático.

\textbf{P: ¿Cómo garantizan la seguridad de los datos sensibles de los socios?} \\
\textbf{R:} Trabajamos con estándares de la industria, encriptando contraseñas, utilizando tokens JWT de corta duración para el manejo de sesiones y asegurando las comunicaciones mediante protocolos TLS. La seguridad no es un agregado, es la base de nuestro desarrollo.

\subsubsection{Axel Zielonka - UX/UI, Frontend, Plataforma Web y App Móvil}

\textbf{P: ¿Cuál es tu rol en el equipo?} \\
\textbf{R:} Soy el puente entre el sistema y el usuario. Me encargo del diseño UX/UI y del desarrollo completo del ecosistema cliente, lo que incluye tanto el panel web para el personal administrativo como la aplicación móvil que usarán los socios en su día a día.

\textbf{P: ¿Qué priorizaste a la hora de diseñar la interfaz?} \\
\textbf{R:} La accesibilidad y la reducción de clics. Los empleados de los clubes suelen rechazar sistemas complejos, por lo que el panel web tiene una curva de aprendizaje mínima. Por el lado de la app móvil, el diseño está centrado en la inmediatez: el socio debe poder mostrar su QR de acceso o pagar su cuota en menos de tres interacciones.

\newpage

\subsection{Preguntas frecuentes específicas por rol}

X.
